\chapter{Antecedentes y Estado del Arte}

\section{Trabajos previos de referencia}

El presente estudio toma como modelo base principal el enfoque propuesto por Moqa \textit{et al.} \cite{moqa2022assessing}, quienes reprodujeron el estudio establecido por Ting \textit{et al.} \cite{ting2010multi} ampliándolo mediante la aplicación de MaOEAs. En su investigación, utilizaron cuatro funciones objetivo para el problema de la selección de etiquetas SNP: la minimización del tamaño del panel genético (compacidad), la maximización de la tolerancia frente a fallos de lectura (tolerancia), la maximización de la distancia media entre secuencias (disimilaritud) y la minimización de la varianza de dicha distancia (balance). Los propios autores justificaron que los MOEAs, como NSGA-II \cite{deb2002fast} y SPEA2 \cite{zitzler2001spea2}, experimentaban una pérdida significativa de presión selectiva al enfrentarse a este elevado número de objetivos, lo que motivó el estudio de MaOEAs, concretamente NSGA-III \cite{deb2013evolutionary} y MOEA/D \cite{zhang2007moea}.

Asimismo, el estudio analizó el rendimiento de estos algoritmos combinándolos con estrategias de inicialización poblacional aleatoria y voraz. El trabajo arrojó cuatro conclusiones fundamentales tras evaluar los algoritmos bajo diversas métricas de calidad (\textit{Range, SumMin, MinSum, MaxToleranceRate, AvgToleranceRate} y \textit{AvgHammingDistance}): en primer lugar, demostró que MOEA/D ofrecía resultados superiores en la mayoría de estas métricas; en segundo lugar, determinó que NSGA-III era el algoritmo superior en la métrica \textit{Max Tolerance Rate}; en tercer lugar, evidenció que SPEA2 superaba al resto en la métrica \textit{Avg Hamming Distance}; y finalmente, reafirmó las ventajas de utilizar la inicialización voraz frente a la aleatoria. El diseño experimental, la formulación de funciones objetivo y el conjunto algorítmico evaluado en dicho trabajo constituyen los cimientos sobre los que se construye este proyecto.

Por otro lado, resulta necesario referenciar el trabajo fundacional de Ting \textit{et al.} \cite{ting2010multi}. Este trabajo fue pionero en formular el problema biológico basándose de forma explícitamente en los cuatro objetivos previamente mencionados, diseñando además la estrategia de inicialización voraz y evaluando las soluciones utilizando el mismo conjunto de métricas que posteriormente heredaría Moqa \textit{et al.} \cite{moqa2022assessing}. En su momento, la investigación evaluó el rendimiento de las soluciones utilizando los algoritmos NSGA-II \cite{deb2002fast} y MSOPS \cite{hughes2003msops}. Las conclusiones de este estudio fueron las siguientes: a nivel algorítmico, demostró que MSOPS superaba a NSGA-II, puesto que este último sufría un exceso de soluciones no dominadas al enfrentarse a cuatro funciones objetivo.

Además, determinó que la inicialización voraz reducía el tiempo de convergencia y mejoraba notablemente a ambos algoritmos, permitiéndoles acercarse a los valores óptimos. En concreto, la combinación de MSOPS con inicialización voraz logró los mejores resultados. A nivel general, el estudio concluyó que la formulación multiobjetivo aportaba una gran flexibilidad para seleccionar diferentes conjuntos de etiquetas SNP, constatando que añadir un número muy pequeño de SNPs adicionales al panel proporcionaba una tolerancia y un balance suficientes para la mayoría de los proyectos.

\section{Problemática detectada en la literatura}

A pesar de que el estudio de Moqa \textit{et al.} \cite{moqa2022assessing} sienta las bases metodológicas para abordar el problema de la selección de etiquetas SNP mediante la evaluación conjunta de MOEA y MaOEA, su estudio presenta una serie de inconsistencias en el reporte de la información. Estas omisiones comprometen la reproducibilidad exacta de sus experimentos y la verificación independiente de sus resultados:

\begin{itemize}
    \item \textbf{Alteración no declarada de las funciones objetivo:} Se observa una fuerte discrepancia entre la formulación matemática teórica de los objetivos y los frentes de Pareto reportados en las gráficas. Concretamente, los rangos de valores observados para la disimilitud ($f_3 \in [0, 0.5]$) y el balance ($f_4 \in [0, 0.04]$) son numéricamente incompatibles con los límites teóricos derivados de sus ecuaciones. Esta incongruencia sugiere que los objetivos han sido alterados matemáticamente, sin que dicha modificación haya sido documentada en la metodología del artículo. Adicionalmente, tampoco detalla la transformación a minimización de los objetivos a maximizar, si es que la ha habido.
    \item \textbf{Ausencia de resultados tabulares:} El estudio omite presentar tablas numéricas con las medias y desviaciones estándar finales alcanzadas por los algoritmos tras las ejecuciones independientes. En lugar de la media, utilizan el rango de convergencia a lo largo de las generaciones, lo cual no aporta información relevante en este dominio, ya que el rendimiento real se evalúa mediante los valores absolutos finales alcanzados. La única forma de aproximar estos valores y establecer un ránking es mediante la inspección visual de las gráficas de convergencia, lo que impide conocer con exactitud el resultado cuantitativo final de las configuraciones evaluadas.
    \item \textbf{Escasez de información paramétrica:} El estudio carece de información estructural indispensable para las arquitecturas evaluadas. Para MOEA/D \cite{zhang2007moea}, no se declara la función de agregación empleada (como suma ponderada o \textit{Tchebycheff}). Adicionalmente, para algoritmos como MOEA/D o NSGA-III \cite{deb2013evolutionary} que emplean vectores de referencia, hay una discrepancia en el tamaño de población reportado en el estudio (200 individuos) con la generación de vectores de referencia mediante el método de Das y Dennis \cite{das1998normal} para cuatro funciones objetivo, ya que la fórmula combinatoria de dicho método no permite generar 200 vectores de referencia uniformes en un espacio de cuatro dimensiones.
    \item \textbf{Lógica de la inicialización voraz:} El estudio tiene una falta de profundidad en la descripción de la inicialización voraz, lo que imposibilita replicar la estrategia de inicialización utilizada.
    \item \textbf{Carencia de métricas de rendimiento globales:} La evaluación del estudio se restringe al análisis de los valores brutos de cada objetivo de forma aislada. No emplea en ningún momento métricas de rendimiento globales propias del estado del arte multiobjetivo, que evalúen simultáneamente tanto la convergencia como la diversidad de los frentes de Pareto descubiertos (tales como el hipervolumen o IGD+).
    \item \textbf{Falta de información sobre la normalización:} En los problemas de optimización de muchos objetivos (MaOPs) donde las funciones presentan escalas de magnitud numéricamente dispares (como ocurre con la compacidad respecto al balance), es un requisito algorítmico imprescindible normalizar el espacio de objetivos para que la evaluación sea matemáticamente justa. Ninguno de los dos trabajos declara ni documenta qué estrategia de normalización se ha empleado en el cálculo de sus resultados empíricos.
    \item \textbf{Falta de accesibilidad al código fuente:} El estudio no proporciona acceso al código \textit{software} empleado para llevar a cabo los experimentos. Esta restricción dificulta la auditoría de la implementación técnica y la verificación de los problemas encontrados.
\end{itemize}
    
Por otro lado, el trabajo fundacional de Ting \textit{et al.} \cite{ting2010multi} tiene carencias similares a las de Moqa \textit{et al.}, como la falta de métricas de rendimiento global y falta de información sobre la normalización, aunque algunas de ellas puedan ser solucionadas accediendo a su código fuente, otras siguen siendo una incógnita.

\begin{itemize}
    \item \textbf{Discrepancia algorítmica no documentada:} En el código fuente original del estudio se evalúa la disimilitud ($f_3$) y el balance ($f_4$) como métricas proporcionales al número de SNPs seleccionados ($f_1$). Concretamente, el código divide la disimilitud bruta entre $f_1$ y el balance entre $f_1^2$. Esta transformación matemática, que convierte la distancia de Hamming absoluta en una métrica de "eficiencia por SNP", se omite por completo en la formulación de las ecuaciones del artículo original y carece de cualquier justificación teórica que respalde su uso.
    \item \textbf{Limitaciones del código fuente proporcionado:} El \textit{software} publicado por Ting \textit{et al.} no contiene toda la información necesaria para replicar la evaluación de su estudio. Su ejecución finaliza en el momento exacto en el que se generan los frentes de Pareto de las soluciones, omitiendo por completo los módulos de evaluación encargados de calcular las métricas finales reportadas en su artículo.
\end{itemize}

Además de las problemáticas expuestas, ninguno de ambos estudios realiza una evaluación rigurosa de sus resultados:

\begin{itemize}
    \item \textbf{Ausencia de validación estadística y escasez de ejecuciones:} Ni el trabajo fundacional de Ting \textit{et al.} ni su posterior extensión por Moqa \textit{et al.} emplean pruebas de significancia estadística (tales como \textit{Kruskal-Wallis}, \textit{Friedman} o \textit{Wilcoxon}) para certificar el rendimiento de los EAs. Las conclusiones sobre la superioridad de una configuración algorítmica frente a otra se fundamentan exclusivamente en la comparación directa de medias tras un número extremadamente reducido de ejecuciones independientes (cinco ejecuciones en el caso de Ting y Moqa). En el ámbito de la computación evolutiva, un tamaño muestral tan pequeño resulta insuficiente para mitigar el ruido estocástico de los algoritmos y carece de la potencia estadística necesaria para validar de forma rigurosa las diferencias de rendimiento, contraviniendo las directrices metodológicas estándar de la literatura que recomiendan realizar al menos 30 ejecuciones independientes \cite{granado2020multiobjective} y otros autores realizan incluso 50 ejecuciones independientes \cite{derrac2011practical}.
\end{itemize}

\section{Contribución de este trabajo}

Este proyecto plantea una implementación del problema de la selección de etiquetas SNP mediante MOEA y MaOEA. La contribución de esta investigación radica en los siguientes puntos:

\begin{itemize}
    \item \textbf{4 funciones objetivo, 2 experimentos:} El estudio plantea dos experimentos, basados en la formulación de las funciones objetivo heredadas de Ting \textit{et al.} \cite{ting2010multi} y Moqa \textit{et al.} \cite{moqa2022assessing}. El primer experimento formula los objetivos de manera absoluta, calculando la disimilitud y el balance sin dividirlos por el número de SNPs seleccionados (compacidad). En contraste, el segundo experimento reproduce la metodología de la literatura, dividiendo dichos objetivos por la compacidad.
    \item \textbf{Integración del marco de optimización Pymoo:} El proyecto moderniza la comparativa de diferentes configuraciones mediante la adopción de la librería estandarizada \texttt{pymoo} \cite{pymoo}, marcando su primera aplicación documentada para este problema genético. El uso de esta librería garantiza implementaciones algorítmicas auditadas de código abierto y computacionalmente estables.
    \item \textbf{Ampliación de los experimentos:} El conjunto de algoritmos evaluado se expande frente a la literatura previa, introduciendo además diferentes operadores de cruce y una nueva estrategia de inicialización puramente voraz, sin componentes aleatorios (\texttt{greedy\_multi}). Además, se amplía la batería de métricas incorporando indicadores de rendimiento individuales (como el hipervolumen \cite{zitzler1999multiobjective}, GD+ e IGD+ \cite{ishibuchi2015igd}) y supra-métricas globales (como el \textit{Average Rank} \cite{demsar2006statistical}).
    \item \textbf{Transparencia de diseño y parametrización:} Los parámetros internos de algoritmos como MOEA/D se encuentran documentados explícitamente, así como la generación de vectores de referencia para los MaOEA que lo necesitan. Adicionalmente, se detallan las estrategias de normalización aplicadas para calcular los resultados de las métricas.
    \item \textbf{Validación estadística:} Para garantizar la validez estadística del análisis, el proyecto realiza 31 ejecuciones independientes por cada configuración evaluada \cite{granado2020multiobjective}. Tras la obtención de cada métrica, se aplican pruebas no paramétricas de significancia estadística, incluyendo los test de \textit{Kruskal-Wallis} \cite{kruskal1952use} y \textit{Friedman} \cite{friedman1937use}, complementados con análisis \textit{post-hoc} de \textit{Dunn} \cite{dunn1964multiple} y \textit{Nemenyi} \cite{nemenyi1963distribution}, respectivamente. Este enfoque asegura que las conclusiones sobre el rendimiento algorítmico estén fundamentadas matemáticamente y no sean producto del azar.

\end{itemize}